\documentclass[12pt,a4paper]{article}
\usepackage[utf8]{inputenc}
\usepackage{ctex}
\usepackage{amsmath}
\usepackage{amssymb}
\usepackage{geometry}
\usepackage{bm}

\geometry{left=2.5cm,right=2.5cm,top=2.5cm,bottom=2.5cm}

% --- 新增的水印所需宏包 ---
\usepackage{graphicx}
\usepackage{xcolor}
\usepackage{eso-pic}

\geometry{left=2cm,right=2cm,top=2.5cm,bottom=2.5cm}

\title{\textbf{AP Calculus AB Practice Test }}
\date{\today}

% --- 用户自定义水印代码开始 ---
\newcommand{\WMText}{贺昱超学习资料}
\newcommand{\WMScale}{2.28}
\newcommand{\WMAngle}{45}
\colorlet{WMColor}{black!8}

\AddToShipoutPictureBG{%
  \AtPageLowerLeft{%
    \put(\LenToUnit{0.66\paperwidth},\LenToUnit{0.70\paperheight}){%
      \rotatebox{\WMAngle}{\scalebox{\WMScale}{\textcolor{WMColor}{\WMText}}}%
    }%
    \put(\LenToUnit{0.33\paperwidth},\LenToUnit{0.50\paperheight}){%
      \rotatebox{\WMAngle}{\scalebox{\WMScale}{\textcolor{WMColor}{\WMText}}}%
    }%
    \put(\LenToUnit{0.66\paperwidth},\LenToUnit{0.30\paperheight}){%
      \rotatebox{\WMAngle}{\scalebox{\WMScale}{\textcolor{WMColor}{\WMText}}}%
    }%
  }%
}
% --- 用户自定义水印代码结束 ---

\title{\textbf{微积分实际工程应用}}

\date{}

\begin{document}
\maketitle
\begin{abstract}
现代 3D 视觉技术（如增强现实（AR）和虚拟现实（VR）的飞速发展，其底层驱动力源自严密的数学理论。主要涉及微积分与高等数学知识在三维空间计算中的核心应用。
如多元微积分（如偏导数、链式法则、一阶泰勒展开）如何构建雅可比矩阵，从而实现连续空间的可微渲染与非线性状态估计；同时，展示了线性代数与群论概念（如奇异值分解 SVD、李群与李代数、舒尔补消元法）在
解决海量参数优化与高维矩阵降维求解中的“工程化魔法”。
\end{abstract}


\vspace{0.5cm}
\hrule
\vspace{0.5cm}

\section{第一阶段：传感器输入 (Sensor Input) —— 睁开眼睛}
这是整个流水线的源头，负责感知现实世界。
\begin{itemize}
    \item \textbf{动作：} 相机以固定的帧率（如 30 FPS）拍摄现实世界，将光信号转化为数字图像序列。
    \item \textbf{数据形态：} 一张张纯粹的 2D 像素矩阵（RGB 图像或灰度图像）。
    \item \textbf{说明：} 如果是双目（Stereo）或深度（RGB-D）相机，此步骤会直接获取初步的深度信息；若是单目相机（Monocular），则需依靠后续算法估算深度（本文以单目为例）。
\end{itemize}

\section{第二阶段：前端视觉里程计 (Visual Odometry) —— 认路与追踪}
前端类似人类的“小脑”，负责处理极短时间内的快速反应，计算相邻两帧之间相机发生的相对运动。

\subsection{步骤 1：特征提取 (Feature Extraction)}
机器无法逐个处理所有像素，必须找出图像中最具辨识度的“特征点”（如桌角、门框）。
\begin{itemize}
    \item \textbf{提取关键点 (Keypoint)：} 寻找图像中梯度变化剧烈的位置（即二维像素坐标 $u, v$）。
    \item \textbf{计算描述子 (Descriptor)：} 考察关键点周围的像素块特征，生成一串数字指纹（如 ORB 的 256 位二进制串），为后续匹配提供依据。
\end{itemize}

\subsection{步骤 2：特征匹配 (Feature Matching)}
当相机移动拍下第 2 帧图像时，机器在两帧的特征点之间进行匹配（对比描述子的汉明距离）。
\begin{itemize}
    \item \textbf{剔除杂音：} 现实中常存在相似纹理干扰导致误匹配。系统采用 \textbf{RANSAC（随机采样一致性）} 算法，剔除不符合物理运动规律的错配点，留下高质量匹配对。
\end{itemize}

\subsection{步骤 3：位姿估计 (Pose Estimation)}
获取匹配的像素点对后，前端根据已知信息的不同，采用三套核心算法求解相机的运动：
\begin{itemize}
    \item \textbf{2D-2D（初始化阶段）：} 系统刚启动且完全无 3D 深度时。使用对极几何（本质矩阵 $E$ 或单应矩阵 $H$），通过八点法解出相机的旋转 $R$ 和平移 $t$，接着用三角测量算出这些点的初始 3D 坐标。
    \item \textbf{3D-2D（正常追踪阶段）：} 上一帧已建好部分 3D 路标点，当前帧仅提取 2D 像素点。系统使用 \textbf{PnP (Perspective-n-Point)} 算法，直接解出当前相机位姿。
    \item \textbf{3D-3D（深度相机专属）：} 两帧均具有 3D 坐标，使用 \textbf{ICP (Iterative Closest Point)} 算法进行匹配。
\end{itemize}

\vspace{0.3cm}
\textbf{【核心数学推导】：2D-2D 前端初始化过程}

\subsubsection{像素坐标归一化}
设图像上的像素坐标为 $\mathbf{p} = [u, v, 1]^T$，相机的内参矩阵为 $K$。为了消除相机内部物理参数的干扰，将像素坐标映射到焦距为 1 的归一化相机坐标平面上，记为 $\mathbf{x}$：
\begin{equation}
    \mathbf{x} = K^{-1} \mathbf{p} = \begin{bmatrix} x \\ y \\ 1 \end{bmatrix}
\end{equation}

\subsubsection{对极约束与八点法 (Epipolar Constraint)}
设空间点 $P$ 在两帧图像上的归一化坐标分别为 $\mathbf{x}_1$ 和 $\mathbf{x}_2$，满足几何刚体变换 $\mathbf{x}_2 \simeq R \mathbf{x}_1 + t$。等式两边同时左乘 $t^{\wedge}$（平移向量的反对称矩阵），再左乘 $\mathbf{x}_2^T$，得到基础的对极约束方程：
\begin{equation}
    \mathbf{x}_2^T t^{\wedge} R \mathbf{x}_1 = 0
\end{equation}
令\textbf{本质矩阵 (Essential Matrix)} 为 $E = t^{\wedge} R$。将 $3 \times 3$ 的矩阵 $E$ 展开为 9 维列向量 $\mathbf{e}$，收集至少 8 对匹配点构造超定线性方程组 $A \mathbf{e} = 0$。通过对 $A$ 进行 \textbf{SVD 分解}，取最小奇异值对应的右奇异向量即可求出 $E$。

\subsubsection{从本质矩阵恢复 $R$ 和 $t$}
对求得的 $E$ 再次进行 SVD 分解：$E = U \Sigma V^T$。引入绕 Z 轴旋转 $\pm \frac{\pi}{2}$ 的矩阵 $W$。平移向量和旋转矩阵可由下式求出（结合深度为正的物理先验排查出唯一真解）：
\begin{equation}
    t^{\wedge} = U W (\text{或 } W^T) \Sigma U^T, \quad R = U W^T (\text{或 } W) V^T
\end{equation}

\subsubsection{直接线性变换三角测量 (DLT Triangulation)}
已知相机位姿后，计算特征点的 3D 空间坐标 $\mathbf{X} = [X, Y, Z, 1]^T$。设两相机的投影矩阵为 $P_1 = K [I | \mathbf{0}]$ 和 $P_2 = K [R | t]$。根据投影模型 $\lambda_1 \mathbf{x}_1 = P_1 \mathbf{X}$ 及 $\lambda_2 \mathbf{x}_2 = P_2 \mathbf{X}$，利用向量叉乘性质消去未知的深度因子 $\lambda$：
\begin{equation}
    \mathbf{x}_1^{\wedge} (P_1 \mathbf{X}) = 0, \quad \mathbf{x}_2^{\wedge} (P_2 \mathbf{X}) = 0
\end{equation}
联立得到形如 $H \mathbf{X} = 0$ 的方程组，再次利用 SVD 分解求得初始 3D 坐标 $\mathbf{X}$。至此，前端给出了一个含有累积误差的相机轨迹和初步的 3D 点云。

\vspace{0.5cm}
\hrule
\vspace{0.5cm}

\section{第三阶段：后端非线性优化 (Back-end Optimization) —— 精雕细琢}
前端只管相邻两帧，目光短浅，误差会随着时间不断累积。后端就像“大脑”，负责全局统筹，消除累积误差。

\subsection{步骤 1：构建关键帧 (Keyframes)}
如果把每一帧（每秒 30 帧）都送给后端优化，电脑会瞬间卡死。所以系统会挑选具有代表性的帧（比如相机移动了一定距离、或者看到了足够多的新特征点）作为\textbf{关键帧}，把冗余的普通帧丢弃。

\subsection{步骤 2：光束法平差 (Bundle Adjustment, BA)}
将一段时间内的所有关键帧位姿和 3D 路标点全部扔进一个巨大的优化方程中，经过 BA 优化后，相机轨迹和 3D 地图就像被“梳理”过一样，变得异常平滑且精确。

\vspace{0.3cm}
\textbf{【核心数学推导】：后端 BA 优化过程}

\subsubsection{后端状态表达：李群与李代数}
相机位姿变换矩阵 $T \in SE(3)$ 不满足加法封闭性。因此引入李代数 $\xi = [\rho, \phi]^T \in \mathfrak{se}(3)$，通过指数映射转换 $T = \exp(\xi^{\wedge})$。
为优化求导，采用\textbf{左扰动模型} (Left Perturbation Model)。给当前位姿 $T$ 施加微小左扰动 $\Delta T = \exp(\delta \xi^{\wedge})$：
\begin{equation}
    T_{new} = \exp(\delta \xi^{\wedge}) T
\end{equation}

\subsubsection{重投影误差与雅可比矩阵推导}
设路标点在世界坐标系下为 $P_w$，投影到像素平面的理论观测值为函数 $h(\xi, P_w)$，实际特征匹配的像素观测值为 $\mathbf{z}$。定义重投影误差 $\mathbf{e}(\xi) = \mathbf{z} - h(\xi, P_w)$。
利用链式法则求误差 $\mathbf{e}$ 对位姿扰动 $\delta \xi$ 的偏导数（雅可比矩阵 $J$），设 $P_c = [X_c, Y_c, Z_c]^T = T P_w$：
\begin{equation}
    J = \frac{\partial \mathbf{e}}{\partial \delta \xi} = \frac{\partial \mathbf{e}}{\partial P_c} \frac{\partial P_c}{\partial \delta \xi}
\end{equation}

第一项（误差对相机坐标系点 $P_c$ 的偏导，$2 \times 3$ 矩阵）：
\begin{equation}
    \frac{\partial \mathbf{e}}{\partial P_c} = - 
    \begin{bmatrix} 
        \frac{f_x}{Z_c} & 0 & -\frac{f_x X_c}{Z_c^2} \\ 
        0 & \frac{f_y}{Z_c} & -\frac{f_y Y_c}{Z_c^2} 
    \end{bmatrix}
\end{equation}

第二项（一阶泰勒展开下，变换点对扰动 $\delta \xi$ 的偏导，$\mathfrak{se}(3)$ 的形式，$3 \times 6$ 矩阵）：
\begin{equation}
    \frac{\partial P_c}{\partial \delta \xi} = 
    \begin{bmatrix} 
        I_{3 \times 3} & -P_c^{\wedge} 
    \end{bmatrix} =
    \begin{bmatrix}
        1 & 0 & 0 & 0 & Z_c & -Y_c \\
        0 & 1 & 0 & -Z_c & 0 & X_c \\
        0 & 0 & 1 & Y_c & -X_c & 0
    \end{bmatrix}
\end{equation}

两项相乘，得到最终的 $2 \times 6$ 位姿雅可比矩阵：
\begin{equation}
    J = -
    \begin{bmatrix} 
        \frac{f_x}{Z_c} & 0 & -\frac{f_x X_c}{Z_c^2} & -\frac{f_x X_c Y_c}{Z_c^2} & f_x + \frac{f_x X_c^2}{Z_c^2} & -\frac{f_x Y_c}{Z_c} \\ 
        0 & \frac{f_y}{Z_c} & -\frac{f_y Y_c}{Z_c^2} & -f_y - \frac{f_y Y_c^2}{Z_c^2} & \frac{f_y X_c Y_c}{Z_c^2} & \frac{f_y X_c}{Z_c} 
    \end{bmatrix}
\end{equation}

\subsubsection{非线性优化求解与舒尔补 (Schur Complement)}
最小化误差平方和，构造高斯-牛顿法 (Gauss-Newton) 的增量方程：$H \delta \mathbf{x} = - \mathbf{b}$。
利用海森矩阵 $H$ 的稀疏箭头结构，将其分块为位姿块 (c) 和路标点块 (p)：
\begin{equation}
    \begin{bmatrix} 
        H_{cc} & H_{cp} \\ 
        H_{pc} & H_{pp} 
    \end{bmatrix} 
    \begin{bmatrix} 
        \delta \mathbf{x}_c \\ 
        \delta \mathbf{x}_p 
    \end{bmatrix} = 
    \begin{bmatrix} 
        -\mathbf{b}_c \\ 
        -\mathbf{b}_p 
    \end{bmatrix}
\end{equation}
通过舒尔补操作，消去对角矩阵 $H_{pp}$ 对应的海量路标点增量，得到仅含位姿增量 $\delta \mathbf{x}_c$ 的方程，大幅降低计算维度，实现实时求解：
\begin{equation}
    (H_{cc} - H_{cp} H_{pp}^{-1} H_{pc}) \delta \mathbf{x}_c = -\mathbf{b}_c + H_{cp} H_{pp}^{-1} \mathbf{b}_p
\end{equation}

\vspace{0.5cm}
\hrule
\vspace{0.5cm}

\section{第四阶段：回环检测 (Loop Closure) —— 唤醒记忆}
即使有后端优化，只要系统运行时间足够长，误差依然会悄悄累积。回环检测就是 SLAM 的“似曾相识（Déjà vu）”机制。
\begin{itemize}
    \item \textbf{词袋模型 (Bag of Words, BoW)：} 系统会把之前看到过的特征点描述子聚类成一本“字典”。当相机来到一个新地方时，它会查字典，判断“这个地方我以前是不是来过？”
    \item \textbf{全局位姿图优化 (Pose Graph Optimization)：} 如果系统确信自己回到了曾经走过的地方（比如绕了一圈回到了原点），它会把当前的轨迹强行“拉”回原点，并将误差均匀地分摊到整个闭环轨迹的每一个节点上，瞬间消除所有累积漂移！
\end{itemize}

\section{第五阶段：建图 (Mapping) —— 描绘世界}
最后，机器要把算出来的 3D 路标点呈现出来供人类或机器使用。
\begin{itemize}
    \item \textbf{稀疏地图 (Sparse Map)：} 就是前端和后端共同维护的那些零散的 3D 特征点。它看起来像一堆骨架，人类很难看懂这是什么，但对于扫地机器人或无人机用来定位和避障已经足够了。
    \item \textbf{稠密地图 (Dense Map)：} 如果你想用在 AR/VR 或者 3D 重建中，稀疏点云是不够的。此时需要开启稠密建图线程，计算图像上每一个像素的深度，拼接成密密麻麻的点云，甚至最终生成带纹理的三维网格（Mesh）模型。
\end{itemize}

\end{document}